% Parton self-energy corrections on the daughters and the soft exchanges between the two hard vertices,
% drawn with TikZ-Feynman. Companion of collinear_diagrams_qq.tex and collinear_diagrams_gg.tex.
% The vertices are placed by hand, so pdflatex suffices. Same set-up as figures/dps_diagram.tex.
% Crossed circles: the hard vertices H_1 and H_2. Double lines: Wilson lines; in panels (i) and (j) these are
% the soft Wilson lines of the four partons k, kbar = p_1 - k, r and rbar. Compile with pdflatex.
% "mgluon" is the gluon coil of TikZ-Feynman reflected about its path. A line or panel that is the
% mirror image of another one under y -> -y uses it, so that the two are exact mirror images.
\documentclass[border=2pt]{standalone}
\pdfinfoomitdate=1 \pdftrailerid{} \pdfsuppressptexinfo=-1     % the same bytes at every build
\usepackage[compat=1.1.0]{tikz-feynman}
\tikzfeynmanset{warn luatex=false, arrow size=1.15pt,
  mgluon/.style={/tikz/draw=none, /tikz/decoration={name=none},
    /tikz/postaction={/tikz/draw, /tikz/decoration={coil, aspect=0.75, segment length=1.5mm}, /tikz/decorate=true}}}
\tikzset{wilson/.style={double, double distance=1.0pt, line width=0.45pt}}
% the two crossed circles, their Wilson lines, and the panel label
\newcommand{\operators}{%
  \vertex[crossed dot] (o1) at (3.0, 0.8) {};
  \vertex[crossed dot] (o2) at (3.0,-0.8) {};}
\newcommand{\wilsonlines}{%
  \draw[wilson] (o1) -- (3.93, 1.29);
  \draw[wilson] (o2) -- (3.93,-1.29);}
\newcommand{\panel}[1]{\node[anchor=north west] at (-0.5, 1.95) {#1};}
\begin{document}
{\small
\begin{tikzpicture}

% ================================================================ self-energies: quark and antiquark; gluon loop on either gluon daughter
% (a) g -> q qbar: self-energy of the quark
\begin{feynman}
  \vertex (p) at (0.0, 0.0);
  \vertex (s) at (1.0, 0.0);
  \vertex (u) at (1.72, 0.288);
  \vertex (v) at (2.56, 0.624);
  \operators
  \diagram* {
    (p) -- [gluon] (s),
    (o2) -- [fermion] (s) -- [plain] (u) -- [fermion] (v) -- [plain] (o1),
    (u) -- [gluon, half left] (v),
  };
\end{feynman}
\wilsonlines
\panel{(a)~quark}

% (b) g -> q qbar: self-energy of the antiquark
\begin{scope}[xshift=4.6cm]
\begin{feynman}
  \vertex (p) at (0.0, 0.0);
  \vertex (s) at (1.0, 0.0);
  \vertex (u) at (1.72,-0.288);
  \vertex (v) at (2.56,-0.624);
  \operators
  \diagram* {
    (p) -- [gluon] (s),
    (o2) -- [plain] (v) -- [fermion] (u) -- [plain] (s) -- [fermion] (o1),
    (u) -- [mgluon, half right] (v),
  };
\end{feynman}
\wilsonlines
\panel{(b)~antiquark}
\end{scope}

% (c) g -> gg: gluon loop on the daughter entering H_1
\begin{scope}[xshift=9.2cm]
\begin{feynman}
  \vertex (p) at (0.0, 0.0);
  \vertex (s) at (1.0, 0.0);
  \vertex (u) at (1.72, 0.288);
  \vertex (v) at (2.56, 0.624);
  \operators
  \diagram* {
    (p) -- [gluon] (s),
    (s) -- [gluon] (u), (v) -- [gluon] (o1),
    (s) -- [mgluon] (o2),
    (u) -- [gluon, half left] (v) -- [gluon, half left] (u),
  };
\end{feynman}
\wilsonlines
\panel{(c)~gluon loop}
\end{scope}

% (d) g -> gg: gluon loop on the daughter entering H_2
\begin{scope}[xshift=13.8cm]
\begin{feynman}
  \vertex (p) at (0.0, 0.0);
  \vertex (s) at (1.0, 0.0);
  \vertex (u) at (1.72,-0.288);
  \vertex (v) at (2.56,-0.624);
  \operators
  \diagram* {
    (p) -- [gluon] (s),
    (s) -- [mgluon] (u), (v) -- [mgluon] (o2),
    (s) -- [gluon] (o1),
    (u) -- [mgluon, half right] (v) -- [mgluon, half right] (u),
  };
\end{feynman}
\wilsonlines
\panel{(d)~gluon loop}
\end{scope}

% ================================================================ self-energies: ghost and quark loops on either gluon daughter
% (e) g -> gg: ghost loop on the daughter entering H_1
\begin{scope}[yshift=-3.8cm]
\begin{feynman}
  \vertex (p) at (0.0, 0.0);
  \vertex (s) at (1.0, 0.0);
  \vertex (u) at (1.72, 0.288);
  \vertex (v) at (2.56, 0.624);
  \operators
  \diagram* {
    (p) -- [gluon] (s),
    (s) -- [gluon] (u), (v) -- [gluon] (o1),
    (s) -- [mgluon] (o2),
    (u) -- [ghost, with arrow=0.5, half left] (v) -- [ghost, with arrow=0.5, half left] (u),
  };
\end{feynman}
\wilsonlines
\panel{(e)~ghost loop}
\end{scope}

% (f) g -> gg: ghost loop on the daughter entering H_2
\begin{scope}[xshift=4.6cm, yshift=-3.8cm]
\begin{feynman}
  \vertex (p) at (0.0, 0.0);
  \vertex (s) at (1.0, 0.0);
  \vertex (u) at (1.72,-0.288);
  \vertex (v) at (2.56,-0.624);
  \operators
  \diagram* {
    (p) -- [gluon] (s),
    (s) -- [mgluon] (u), (v) -- [mgluon] (o2),
    (s) -- [gluon] (o1),
    (u) -- [ghost, with arrow=0.5, half right] (v) -- [ghost, with arrow=0.5, half right] (u),
  };
\end{feynman}
\wilsonlines
\panel{(f)~ghost loop}
\end{scope}

% (g) g -> gg: quark loop on the daughter entering H_1
\begin{scope}[xshift=9.2cm, yshift=-3.8cm]
\begin{feynman}
  \vertex (p) at (0.0, 0.0);
  \vertex (s) at (1.0, 0.0);
  \vertex (u) at (1.72, 0.288);
  \vertex (v) at (2.56, 0.624);
  \operators
  \diagram* {
    (p) -- [gluon] (s),
    (s) -- [gluon] (u), (v) -- [gluon] (o1),
    (s) -- [mgluon] (o2),
    (u) -- [fermion, half left] (v) -- [fermion, half left] (u),
  };
\end{feynman}
\wilsonlines
\panel{(g)~quark loop}
\end{scope}

% (h) g -> gg: quark loop on the daughter entering H_2
\begin{scope}[xshift=13.8cm, yshift=-3.8cm]
\begin{feynman}
  \vertex (p) at (0.0, 0.0);
  \vertex (s) at (1.0, 0.0);
  \vertex (u) at (1.72,-0.288);
  \vertex (v) at (2.56,-0.624);
  \operators
  \diagram* {
    (p) -- [gluon] (s),
    (s) -- [mgluon] (u), (v) -- [mgluon] (o2),
    (s) -- [gluon] (o1),
    (u) -- [fermion, half right] (v) -- [fermion, half right] (u),
  };
\end{feynman}
\wilsonlines
\panel{(h)~quark loop}
\end{scope}

% ================================================================ soft exchanges between Wilson lines entering different hard vertices
% (i) soft gluon between the Wilson lines of k and rbar
\begin{scope}[xshift=4.6cm, yshift=-7.6cm]
\begin{feynman}
  \vertex[crossed dot] (h1) at (1.75, 1.0) {};
  \vertex[crossed dot] (h2) at (1.75,-1.0) {};
  \vertex (a) at (0.95, 0.61);
  \vertex (b) at (2.55,-0.61);
  \diagram* {
    (a) -- [gluon] (b),
  };
\end{feynman}
\draw[wilson] (0.15, 0.22) -- (h1);   % k
\draw[wilson] (0.15,-0.22) -- (h2);   % kbar = p_1 - k
\draw[wilson] (3.35, 0.22) -- (h1);   % r
\draw[wilson] (3.35,-0.22) -- (h2);   % rbar
\node at (0.50, 0.78) {\(k\)};
\node at (0.50,-0.80) {\(\bar k\)};
\node at (3.00, 0.78) {\(r\)};
\node at (3.00,-0.80) {\(\bar r\)};
\node at (1.75, 1.42) {\(H_1\)};
\node at (1.75,-1.42) {\(H_2\)};
\panel{(i)~soft}
\end{scope}

% (j) soft gluon between the Wilson lines of kbar and r
\begin{scope}[xshift=9.2cm, yshift=-7.6cm]
\begin{feynman}
  \vertex[crossed dot] (h1) at (1.75, 1.0) {};
  \vertex[crossed dot] (h2) at (1.75,-1.0) {};
  \vertex (a) at (0.95,-0.61);
  \vertex (b) at (2.55, 0.61);
  \diagram* {
    (a) -- [mgluon] (b),
  };
\end{feynman}
\draw[wilson] (0.15, 0.22) -- (h1);   % k
\draw[wilson] (0.15,-0.22) -- (h2);   % kbar = p_1 - k
\draw[wilson] (3.35, 0.22) -- (h1);   % r
\draw[wilson] (3.35,-0.22) -- (h2);   % rbar
\panel{(j)~soft}
\end{scope}

\end{tikzpicture}}
\end{document}
